% Section 8: numerical experiments. Sources (all under test/bie-bound/):
% litbench-v1, pole-free-v1, rayleigh-order-v1, solver-cost-v1,
% failure-probe-v1 (P1), wood-fallback-v1. Numbers are copied from
% the experiment READMEs; the section heading is in main.tex.

\subsection{Examples and settings}\label{sec:examples}
All examples have $d=1$ and $\rho_{\mathrm e}=1$, so that only the order $m=0$
propagates in the windows used and its Wood line is $k=|\beta|$.
\begin{itemize}
\item \emph{Example 1}: $\beta=0.5$, $L_0=L_\pm=1$. The defect rod, with $\rho=4$,
  is the star
  $(x,y)=0.2\,(1.5\cos t+0.15\cos6t+0.15\cos4t,\ \sin t+0.15\sin6t-0.15\sin4t)$.
  Each periodicity cell contains, at its center, the three-lobed rod
  $r(t)=c\,(1+a\cos3t)$ with $c=0.2(1+a^2/2)^{-1/2}$ and $\rho=17$, with $a=0.12$ on
  the left and $a=0.20$ on the right.
\item \emph{Example 2}: as Example 1, but with an empty defect cell, $L_+=0.85$, and
  the rod of $S_+$ rotated by $0.13$.
\item \emph{W3}: the guide with three missing rows of \cite{BottenEtAl2006}, a
  hexagonal array of circular rods of radius $0.2$ and refractive index $3$ in
  air, TM polarization, row spacing $h=\sqrt3/2$, half-strip period $2h$.
\end{itemize}
The smallest gap widths are $w_{\min}=0.4685$, $0.4289$ and $h-0.4=0.4660$. A
grid label $N/M$ gives the number of nodes on the defect rod and the
Rayleigh half-width; the rods of the half-strips have $N/2$ nodes. Unless stated otherwise the
two grids are $128/6$ and $192/12$, and a root is accepted only if it is found
on both (\cref{sec:search}). Timings are wall-clock times on one machine,
including the start and exit of MATLAB.

\subsection{Literature benchmark}\label{sec:litbench}
For W3 at $d/\lambda=0.32$ the solver gives $\alpha_0d=\beta d=1.6754070619$,
which agrees with the printed value $1.67541$ of \cite{BottenEtAl2006} in all
six digits. Conversely, at the printed $\beta d=1.67541$ it gives
$d/\lambda=0.32000027$. The second mode at $\beta d=0.2$ gives
$d/\lambda=0.33252418$, consistent with the published dispersion figure. For
$N\in\{64,128,256\}$ and $M\in\{12,24\}$ the six computed values of each mode
agree to at least $15$ significant digits (\texttt{litbench-v1}).

\subsection{Pole-free closure}\label{sec:exp-pole}
\Cref{tab:pole-free} compares the Dirichlet-to-Neumann closure
\eqref{eq:port-matrix} with the bordered closure \eqref{eq:bordered} on
regular and moved artificial boundaries of Example 1. Over nine cases at two resolutions,
the bordered search returned every expected root, accepted no pole, and
refused no window. The root counts of $18$ independent dense scans of the
projection test \eqref{eq:projection} agreed with the search. In the regular
cases the two closures agree to $8.4\times10^{-15}$.

\begin{table}[htbp]
\centering\small
\caption{Example 1 near $k_*=1.558437894182088$ with the right artificial
boundary moved inside its gap strip. ``Offset'' places a Dirichlet pole $0.0022$ below
the mode; ``aligned'' moves the pole onto the mode; ``node'' puts the root on
a Chebyshev node. The last column is the wall jump of the accepted field.
Source: \texttt{pole-free-v1}.}
\label{tab:pole-free}
\begin{tabular}{llll}
\toprule
case, grid & DtN closure & bordered closure & wall jump\\
\midrule
reference, 128/6 & $1.558437894182088$ & $1.5584378941820876$ & $4.3\times10^{-6}$\\
offset, 128/6 & refused & $1.5584378941820884$ & $4.3\times10^{-5}$\\
offset, 192/12 & $1.55625026$ (pole) & $1.5584378941820887$ & $1.5\times10^{-7}$\\
aligned, 128/6 & refused & $1.5584378941820889$ & $4.5\times10^{-5}$\\
aligned, 192/12 & refused & $1.5584378941820887$ & $1.6\times10^{-7}$\\
node, 192/12 & no root & $1.5584378941820889$ & $5.3\times10^{-11}$\\
\bottomrule
\end{tabular}
\end{table}

The Dirichlet-to-Neumann closure accepted the pole at $1.55625026$ because
its singular-value ratio there was $4.5\times10^{-16}$, the largest singular
value having grown to about $2\times10^{13}$. The physical test of
\cref{sec:search} also rejects this frequency: the field of the old null
vector has wall jumps $0.052$ and $0.061$ on the two grids, which fails both
the absolute and the refinement gate. The bordered matrix has no small
singular value there ($\sigma_{\min}=1.75\times10^{-3}$, against about
$10^{-14}$ at roots), so it never proposes this frequency. The field test is a
second line of defense; it is the structure of \eqref{eq:bordered} that
removes the pole.

\subsection{Rayleigh order}\label{sec:exp-order}
\Cref{tab:order} gives $|k_M-k_{24}|$ at fixed $N$. Fitting
$\log|k_M-k_{24}|\approx c-2\pi Mw_{\mathrm{eff}}$ over the orders above the
rounding floor gives $w_{\mathrm{eff}}=0.738$, $0.698$, $0.928$ for Example 1,
Example 2 and W3. These are larger than $w_{\min}$ in every case, which
\cref{thm:root-error} allows and the continuation remark of
\cref{sec:truncation} explains in part. Moving the right artificial boundary of
Example 1 by $-0.04$, $+0.06$ and $+0.16$ changes the roots by at most
$1.6\times10^{-15}$ and $w_{\mathrm{eff}}$ by less than $10^{-4}$, as predicted by
the paragraph on the positions of the artificial boundaries in \cref{sec:truncation}. In the common gaps of
all three examples the computed Floquet eigenvalue problems had $K$ multipliers inside the
unit circle, with flux defects $\|F\|$ of about $10^{-15}$
(\texttt{rayleigh-order-v1}).

\begin{table}[htbp]
\centering\small
\caption{$|k_M-k_{24}|$ at $N=128$. Source: \texttt{rayleigh-order-v1}.}
\label{tab:order}
\begin{tabular}{lccccc}
\toprule
 & $M=1$ & $2$ & $3$ & $4$ & $5$\\
\midrule
Example 1 & $9.1\times10^{-7}$ & $1.5\times10^{-9}$ & $4.9\times10^{-11}$ & $5.5\times10^{-13}$ & $2.2\times10^{-15}$\\
Example 2 & $1.0\times10^{-6}$ & $1.4\times10^{-9}$ & $1.3\times10^{-11}$ & $6.3\times10^{-13}$ & $1.4\times10^{-14}$\\
W3 & $5.2\times10^{-6}$ & $1.3\times10^{-8}$ & $3.9\times10^{-11}$ & $1.3\times10^{-13}$ & $2.2\times10^{-15}$\\
\bottomrule
\end{tabular}
\end{table}

\subsection{Cost against high-order finite elements}\label{sec:exp-cost}
The task is that of \cref{sec:task} for Example 1: return all guided modes in
the window $[1.451,2.249]$ at $\beta=0.5$, starting from a fresh MATLAB process
with no cached data, on $11$ threads. The comparison method is the finite
element method with $P_4$ elements on a mesh with $g$ background subdivisions
per cell side and the Dirichlet-to-Neumann closure of
\cite{JolyLiFliss2006,Fliss2013}. Each solver's error is the distance of its
root to the reference value $k_{\mathrm{ref}}=1.558437894182088$. All boundary
integral roots with $N\in\{128,192,256,512\}$ and $6\le M\le48$ agree with
$k_{\mathrm{ref}}$ to within $9\times10^{-16}$, and the $P_4$ roots approach it
under mesh refinement, with distances $1.5\times10^{-9}$, $2.9\times10^{-11}$ and
$1.8\times10^{-12}$ for $g=12,24,48$. For the boundary integral solver the error is
therefore a self-convergence estimate, checked independently by the finite
element roots to about $2\times10^{-12}$; it is not a bound, and errors below
about $10^{-12}$ are not resolved.
\Cref{tab:cost} gives the cheapest measured configuration of each method in
each accuracy class.

\begin{table}[htbp]
\centering\small
\caption{Cheapest configuration reaching each error class, Example 1. Time
includes MATLAB start and exit; each configuration was timed once.
The error is $|k-k_{\mathrm{ref}}|$. Source:
\texttt{solver-cost-v1}.}
\label{tab:cost}
\begin{tabular}{lllr}
\toprule
error & BIE: grid, s, GiB & FEM: grid, s, GiB & BIE/FEM time\\
\midrule
$10^{-6}$ & $128/2$, $38.6$, $0.84$ & $P_4$, $g=12$, $16.7$, $1.05$ & $2.31$\\
$10^{-8}$ & $128/2$, $38.6$, $0.84$ & $P_4$, $g=12$, $16.7$, $1.05$ & $2.31$\\
$10^{-10}$ & $128/6$, $38.8$, $0.81$ & $P_4$, $g=24$, $27.5$, $1.29$ & $1.41$\\
$10^{-11}$ & $128/6$, $38.8$, $0.81$ & $P_4$, $g=48$, $97.8$, $1.21$ & $0.40$\\
\bottomrule
\end{tabular}
\end{table}

For errors between $10^{-6}$ and $10^{-10}$ the finite element method is $1.4$
to $2.3$ times faster. The time of the boundary integral solver hardly
depends on the accuracy, about $39$\,s, most of it in $35$ determinant
evaluations; with $M=2$ and $M=6$ it takes $38.6$ and $38.8$\,s. For an error
below $10^{-11}$ it is $2.5$ times faster than the cheapest finite element
configuration that reaches it. On this example the boundary integral
solver is therefore the cheaper method only at the highest accuracy. The
comparison covers one example, one machine and one timing per
configuration.

\subsection{Near Wood points}\label{sec:exp-wood}
The guided branch of Example 1 crosses the Wood line $k=|\beta|$ near
$\beta\approx1.98154$, and that of W3 near $\beta\approx2.97503$
(\texttt{failure-probe-v1}). At $\beta$ values whose root lies at distance
$10^{-4}$ to $10^{-1}$ from $k_W$, a Newton iteration started near the root
converged on both grids. At distance $10^{-4}$ the two grids agree to
$6.4\times10^{-12}$ (Example 1) and $2.7\times10^{-12}$ (W3). The global search behaved worse. A window
containing $k_W$ failed at its first evaluation, at $k=k_W$. One-sided
windows ending $10^{-4}$ from $k_W$ did not finish within $130$\,s (coarse
grid) and $90$\,s (fine grid) in any of eight cases.

The fallback of \cref{sec:wood-fallback} was tested on windows
$[k_W-0.02,k_W+0.02]$ at the $\beta$ values above, eight for Example 1 and
four for W3, on both grids, with a time cap of $150$\,s per search and grid
(\texttt{wood-fallback-v1}). The two Example 1 roots at distance $0.1$ were
searched in separate windows. No evaluation fell inside an excluded band;
in $12{,}498$ cell evaluations the closest one was at distance $\delta$,
up to rounding.

\emph{Completion.} With $\delta=10^{-3}$ every search finished, in both
interpolation variables: $24$ of $24$ cases, with median cumulative determinant times of
$26$\,s for the variable $u=\operatorname{sign}(k-k_W)|k-k_W|^{1/2}$ and $37$\,s for
$k$ itself. At $\delta=10^{-4}$ the outcome depended on the panels. When the
retained side of $k_W$ was one $u$-panel from the start, all $30$ searches of
the branch scan below finished, in $17$ to $38$\,s. In one Example 1 window the
$k$-panels needed $308$ evaluations and $122$\,s, against $56$ evaluations and
$21$\,s in $u$. When a search at $\delta=10^{-2}$ was extended by the two new
strips down to $10^{-4}$, five of the six tested cases in each variable did
not finish within the cap or reached the depth limit. At $\delta=10^{-5}$ no tested search finished. In
these strips the coefficient and off-grid tests kept failing under
subdivision (for example a relative off-grid discrepancy of $3\times10^{-9}$
against the gate $10^{-9}$). This is consistent with the Green error of
\cref{sec:wood-mechanism}, which subdivision cannot remove, but we have not
isolated the cause.

\emph{Accuracy.} \Cref{tab:wood-roots} compares each returned root with a
reference computed with a more accurate Green function: the proxy fit of the
quasi-periodic Green function is replaced by one that prescribes its
exponentially convergent Rayleigh series on the top and bottom of the fitting
box. This reduces the Green error at $t=10^{-4}$ from $1.2\times10^{-10}$ to
$7.1\times10^{-15}$; the boundary integral discretization is unchanged, so the
difference measures only the effect of the Green error. Refining the new fit
further moved the reference roots by at most $1.7\times10^{-12}$. The
difference grows roughly like $1/t$
(fitted slopes $-0.94$ and $-1.11$ on the two sides of $k_W$), and
$t\,|k-k_{\mathrm{ref}}|\le4.1\times10^{-14}$ for Example 1 and
$5.7\times10^{-15}$ for W3 over the tested range. Near $k_W$ the two-grid
difference underestimates this error by a factor of about $16$, as
\cref{sec:wood-mechanism} predicts.

\begin{table}[htbp]
\centering\small
\caption{Roots returned by the fallback on the fine grid, Example 1.
$t=k_*-k_W$; ``reference'' uses the more accurate Green function on the same
discretization. The rows $\pm1.2\times10^{-4}$ are from the branch scan
($\delta=10^{-4}$), the others from searches with $\delta=10^{-3}$ or $10^{-4}$.
Source: \texttt{wood-fallback-v1}.}
\label{tab:wood-roots}
\begin{tabular}{rlll}
\toprule
$t$ & returned $k$ & two-grid difference & difference to reference\\
\midrule
$10^{-1}$ & $1.9117673817423038$ & $2.2\times10^{-16}$ & $1.8\times10^{-13}$\\
$10^{-2}$ & $1.9743542511171162$ & $6.0\times10^{-15}$ & $1.8\times10^{-12}$\\
$10^{-3}$ & $1.9808233223406590$ & $<10^{-16}$ & $7.6\times10^{-12}$\\
$1.2\times10^{-4}$ & $1.9814575417863407$ & $3.8\times10^{-12}$ & $6.3\times10^{-11}$\\
$-1.2\times10^{-4}$ & $1.9816305613218432$ & $4.0\times10^{-12}$ & $6.5\times10^{-11}$\\
$-10^{-3}$ & $1.9822651530044606$ & $7.0\times10^{-14}$ & $8.3\times10^{-12}$\\
$-10^{-2}$ & $1.9887716011465046$ & $<10^{-16}$ & $1.1\times10^{-12}$\\
$-10^{-1}$ & $2.0549852982808585$ & $3.6\times10^{-15}$ & $3.8\times10^{-14}$\\
\bottomrule
\end{tabular}
\end{table}

\emph{Branch check.} For $15$ values of $\beta$ along the Example 1 branch,
with $k_*-k_W$ between $-5\times10^{-4}$ and $5\times10^{-4}$, the fallback with
$\delta=10^{-4}$ returned the mode on both grids at the $10$ values where it lies
outside the excluded band. At the other five, where interpolation from the
neighboring $\beta$ puts it at $|k_*-k_W|\le8\times10^{-5}$, both grids
finished their retained windows with no root, and the band was reported as
not searched. No other root was returned near either edge of the band.

The fallback thus does what \cref{sec:wood-fallback} claims, on two examples:
it never evaluates the cell matrices near $k_W$, it returns the modes outside
the band, and it reports the band. Its accuracy near $k_W$ is limited by the
Green error, not by the interpolation, and the two-grid test does not see this
limit.
